\section{Divisor-defect feedback from the actual residual (NS-75)}
\label{sec:ns75-divisor-feedback}

\paragraph{Question.}
The fixed M\"obius templates in NS-73 do not use the optimized old
coefficients. We now test a correction that does: cancel the next block
of derivative jumps of the actual residual. This yields an exact
arithmetic construction, but removing those local jumps is not a
norm-descent theorem.

Write the complete residual as
$R_N=\tau-\sum_{n\leq N}c_n\sigma_n$, with the actual optimized $c_n$,
and put
\[
 \eta_c=\sum_{n\leq N}\frac{c_n}{n},\qquad
 r_N(m)=\mathbf1_{m=1}+\sum_{\substack{n\mid m\\n\leq N}}c_n.
\]
The following statements also hold for arbitrary fixed real old coefficients.

\begin{proposition}[Complete divisor-log cells and derivative jumps]
\label{prop:ns75-divisor-cells}
For every $x>0$,
\begin{equation}
 R_N(x)=-\frac{\eta_c}{x}
       +\sum_{m\leq1/x}r_N(m)\log\frac1{mx}.
 \label{eq:ns75-divisor-cells}
\end{equation}
At $x=1/m$ the function is continuous and the right-minus-left jump of
its first derivative is $m r_N(m)$. In distributions on $(0,\infty)$,
\begin{equation}
 R_N''=\frac1{x^2}\sum_{m\leq1/x}r_N(m)
           -\frac{2\eta_c}{x^3}
           +\sum_{m\geq1}m r_N(m)\delta_{1/m}.
 \label{eq:ns75-jump-measure}
\end{equation}
The atomic sum is locally finite away from zero. The complete exterior
tail is $R_N(x)=-\eta_c/x$ for $x>1$.
\end{proposition}
\begin{proof}
The exact identity
$A(z)=z-\sum_{k\leq z}\log(z/k)$ gives
\[
 \sum_{n\leq N}c_n\sigma_n(x)
 =\frac{\eta_c}{x}-\sum_{m\leq1/x}
       \left(\sum_{\substack{n\mid m\\n\leq N}}c_n\right)
       \log\frac1{mx}.
\]
For each fixed $x$ both sums are finite, so this rearrangement requires
no limiting interchange. The target is exactly the term $m=1$ in the
logarithmic sum. A term entering at $1/m$ has value zero there, proving
continuity. Between endpoints, differentiation gives
$R_N'=\eta_c/x^2-x^{-1}\sum_{m\leq1/x}r_N(m)$.
The stated jump and second-derivative distribution follow, including
the target endpoint $m=1$.
\end{proof}

\paragraph{An exact local cancellation step.}
For $N<n\leq2N$ choose physical new coefficients
\[
 a_n^{(p)}=-r_N(n)\log^p(2N/n),\qquad p=0,1,2,
\]
with the constant taper convention for $p=0$.
With $p=0$, append these coefficients while holding all old coefficients
fixed. For $N<m\leq2N$, every proper divisor of $m$ is at most $N$.
Consequently the updated divisor defect is exactly
$r_N(m)+a_m^{(0)}=0$ throughout this next block. The old defects are
unchanged. Other cells, the regular part of~\eqref{eq:ns75-jump-measure}
and both tails remain in the norm calculation. No smallness follows
from cancellation of the displayed atomic part alone.

If the old coefficients happen to be $c_n=-\mu(n)$, then
$r_N(m)=\mu(m)$ on this next block, so the constant rule specializes
to NS-73's sharp M\"obius template. This is a specialization, not an
assumption about the optimized old coefficients.

For comparison with the optimum we also convert each $a^{(p)}$ to
difference coordinates using~\eqref{eq:ns73-coordinate-change}, apply
the complete old-space complement, and optimize its scalar step.
The exact gain and the best three-direction span gain are then those
of Proposition~\ref{prop:ns73-span}. This old-space reoptimization
can change the local jumps again; we do not claim it preserves the
raw cancellation property. We separately record the complete error
of the held-old, unit-sized step:
\[
 \frac{\|R_N-\sum_{n=N+1}^{2N}a_n^{(p)}\sigma_n\|_2^2}{E_N}
 =1+\frac{(a^{(p)})^TG_{\rm new,new}a^{(p)}
                 -2(g^Tv^{(p)})}{E_N}.
\]
Old-space orthogonality justifies the numerator equality here;
the raw physical Gram in this expression is not replaced by a
projected or model Gram.

\paragraph{Certified finite outcome.}
The complete 256/384-bit replays and doubled-cutoff replay at $N=256$
give the following strict enclosures for the fraction of full block gain.
The single cancellation direction is optimally scaled \emph{after} old-space
projection; its scalar is not constrained positive.
\begin{center}
\begin{tabular}{rccc}
\hline
$N$&Cancellation direction&Best three-direction span&Diagonal curvature\\
\hline
16&(.3110,.3111)&(.5836,.5837)&(.9308,.9310)\\
32&(.4292,.4294)&(.8086,.8087)&(.9040,.9041)\\
64&(.6954,.6956)&(.9063,.9065)&(.9783,.9784)\\
128&(.0029,.0030)&(.7232,.7233)&(.9852,.9854)\\
256&(.0010,.0011)&(.5602,.5603)&(.9763,.9764)\\
\hline
\end{tabular}
\end{center}
At $N=256$ the span retains $56.02048\%$ of full gain, versus
$97.63068\%$ for curvature. The raw held-old cancellation step increases
the complete squared error by a factor in $(1495.54,1495.55)$.
Indeed all fifteen held-old unit steps tested here increase squared error.
The cancellation direction's correlation at $N=256$ is negative, so its
best scalar step reverses the prescribed sign. Merely cancelling derivative
jumps is therefore not a descent principle even on these certified blocks.
This is a finite failure of the proposed unit step and finite underperformance
of its span, not an asymptotic exclusion of every divisor-feedback scheme.
Eleven exact Maxima checks cover the identities and coordinate conversion;
all finite gains retain complete tails, the true old projection and all
three-direction cross terms.
